\section{流形的嵌入}\label{section:ch04-manifold-embedding}

我们已经证明了每一个具有可数基的正则空间可以嵌入到``无穷维''欧氏空间 $\mathbb{R}^{\omega}$ 中。一个自然的问题是在怎样的条件下空间 $X$ 可以被嵌入到某一个有限维欧氏空间 $\mathbb{R}^{N}$。本节我们将回答上述问题。在第 8 章我们研究维数论时，将对这个问题给出一个更具一般性的回答。

\begin{definition}\label{definition:ch04-manifold}
一个 $m$\textbf{-维流形}（$m$-manifold）是指一个具有可数基的 Hausdorff 空间 $X$，它的每一点 $x$ 有一个邻域同胚于 $\mathbb{R}^{m}$ 中的一个开子集。
\end{definition}

$1$-维流形通常称为\textbf{曲线}（curve），$2$-维流形称为\textbf{曲面}（surface）。流形是一类很重要的空间。在微分几何和代数拓扑中有充分的研究。

我们将证明，若 $X$ 是一个紧致流形，则 $X$ 可以嵌入到一个有限维欧氏空间中。去掉紧致性的假设，这个定理也成立，但是证明将变得相当困难。

首先，我们需要介绍某些术语。

若 $\phi\colon X\to\mathbb{R}$，则 $\phi$ 的\textbf{支撑}（support）被定义为集合 $\phi^{-1}(\mathbb{R}-\{\boldsymbol{0}\})$ 的闭包。因此若 $x$ 在 $\phi$ 的支撑之外，则 $x$ 有一个邻域，在这个邻域上 $\phi$ 蜕化。

\begin{definition}\label{definition:ch04-partition-of-unity}
设 $\{U_{1},\cdots,U_{n}\}$ 是空间 $X$ 的一个加标有限开覆盖。连续函数的一个加标族
\begin{align*}
\phi_{i}\colon X\longrightarrow[0,1]\quad\text{对于 }i=1,\cdots,n,
\end{align*}
称为由 $\{U_{i}\}$ 控制的一个\textbf{单位分拆}（partition of unity），如果：
\begin{enumerate}
\item 对于每一个 $i$，$(\phi_{i}\text{ 的支撑})\subset U_{i}$.
\item 对于每一个 $x$，$\sum_{i=1}^{n}\phi_{i}(x)=1$.
\end{enumerate}
\end{definition}

\begin{theorem}[有限单位分拆的存在性]\label{theorem:ch04-partition-of-unity}
设 $\{U_{1},\cdots,U_{n}\}$ 是正规空间 $X$ 的一个有限开覆盖。则存在一个由 $\{U_{i}\}$ 控制的单位分拆。
\end{theorem}

\begin{proof}
第一步。首先证明覆盖 $\{U_{i}\}$ 能``缩小''为 $X$ 的一个开覆盖 $\{V_{1},\cdots,V_{n}\}$，使得对于每一个 $i$ 有 $\overline{V}_{i}\subset U_{i}$。

用归纳法证明。首先注意：集合
\begin{align*}
A=X-(U_{2}\cup\cdots\cup U_{n})
\end{align*}
是 $X$ 的一个闭子集。因为 $\{U_{1},\cdots,U_{n}\}$ 覆盖 $X$，所以集合 $A$ 被包含在开集 $U_{1}$ 中。应用正规性，选取一个包含 $A$ 的开集 $V_{1}$，使得 $\overline{V}_{1}\subset U_{1}$。于是集族 $\{V_{1},U_{2},\cdots,U_{n}\}$ 覆盖 $X$。

一般地，设开集 $V_{1},\cdots,V_{k-1}$ 已经给定，使得集族
\begin{align*}
\{V_{1},\cdots,V_{k-1},U_{k},U_{k+1},\cdots,U_{n}\}
\end{align*}
覆盖 $X$。令
\begin{align*}
A=X-(V_{1}\cup\cdots\cup V_{k-1})-(U_{k+1}\cup\cdots\cup U_{n}).
\end{align*}
则 $A$ 是 $X$ 的一个闭集，并且包含在开集 $U_{k}$ 中。选取 $V_{k}$ 为包含 $A$ 的一个开集，使得 $\overline{V}_{k}\subset U_{k}$。于是 $\{V_{1},\cdots,V_{k-1},V_{k},U_{k+1},\cdots,U_{n}\}$ 覆盖 $X$。经 $n$ 步归纳，我们便可证明结论。

第二步。现在来证明定理。给定 $X$ 的一个开覆盖 $\{U_{1},\cdots,U_{n}\}$，选取 $X$ 的一个开覆盖 $\{V_{1},\cdots,V_{n}\}$，使得对于每一个 $i$ 有 $\overline{V}_{i}\subset U_{i}$。再选取 $X$ 的一个开覆盖 $\{W_{1},\cdots,W_{n}\}$，使得对于每一个 $i$ 有 $\overline{W}_{i}\subset V_{i}$。应用 Urysohn 引理，对于每一个 $i$，选取一个连续函数
\begin{align*}
\psi_{i}\colon X\longrightarrow[0,1]
\end{align*}
使得 $\psi_{i}(\overline{W}_{i})=\{1\}$ 和 $\psi_{i}(X-V_{i})=\{0\}$ 成立。因为 $\psi_{i}^{-1}(\mathbb{R}-\{\boldsymbol{0}\})$ 包含在 $V_{i}$ 中，从而
\begin{align*}
(\psi_{i}\text{ 的支撑})\subset\overline{V}_{i}\subset U_{i}.
\end{align*}
又由于集族 $\{W_{i}\}$ 覆盖 $X$，对于每一个 $x$，$\Psi(x)=\sum_{i=1}^{n}\psi_{i}(x)$ 取正值。因此，对于每一个 $j$，定义
\begin{align*}
\phi_{j}(x)=\frac{\psi_{j}(x)}{\Psi(x)}.
\end{align*}
容易验证函数 $\phi_{1},\cdots,\phi_{n}$ 便是我们所需的单位分拆。
\end{proof}

有一个与单位分拆相似的概念，其中开覆盖和函数族不是有限的，甚至不是可数的。在第 6 章研究仿紧致性时，我们将讨论相关问题。

\begin{theorem}\label{theorem:ch04-manifold-embedding}
若 $X$ 是一个 $m$-维紧致流形，则 $X$ 可以嵌入到 $\mathbb{R}^{N}$ 中，其中 $N$ 是某一个正整数。
\end{theorem}

\begin{proof}
用有限多个开集 $\{U_{1},\cdots,U_{n}\}$ 覆盖 $X$，其中每一个开集都可以嵌入到 $\mathbb{R}^{m}$ 中。对于每一个 $i$，选取一个嵌入 $g_{i}\colon U_{i}\to\mathbb{R}^{m}$。由于 $X$ 是一个紧致的 Hausdorff 空间，所以它是正规的。设 $\phi_{1},\cdots,\phi_{n}$ 是由 $\{U_{i}\}$ 控制的一个单位分拆。令 $A_{i}=\phi_{i}$ 的支撑。对于 $i=1,\cdots,n$，定义函数 $h_{i}\colon X\to\mathbb{R}^{m}$ 为
\begin{align*}
h_{i}(x)=\begin{cases}
\phi_{i}(x)\cdot g_{i}(x) & \text{当 }x\in U_{i},\\
\boldsymbol{0}=(0,\cdots,0) & \text{当 }x\in X-A_{i}.
\end{cases}
\end{align*}
[这里 $\phi_{i}(x)$ 是一个实数 $c$，$g_{i}(x)$ 是 $\mathbb{R}^{m}$ 的一个点 $\boldsymbol{y}=(y_{1},\cdots,y_{m})$。显然，乘积 $c\cdot\boldsymbol{y}$ 是 $\mathbb{R}^{m}$ 中的点 $(cy_{1},\cdots,cy_{m})$。] 由于决定 $h_{i}$ 的两个函数在它们的定义域的交上取值相同，因此 $h_{i}$ 的定义是确切的。由于 $h_{i}$ 在开集 $U_{i}$ 和 $X-A_{i}$ 上的限制都是连续的，所以 $h_{i}$ 也是连续的。

现在定义函数
\begin{align*}
F\colon X\longrightarrow\big(\underbrace{\mathbb{R}\times\cdots\times\mathbb{R}}_{n\text{个}}\times\underbrace{\mathbb{R}^{m}\times\cdots\times\mathbb{R}^{m}}_{n\text{个}}\big)
\end{align*}
为
\begin{align*}
F(x)=(\phi_{1}(x),\cdots,\phi_{n}(x),h_{1}(x),\cdots,h_{n}(x)).
\end{align*}
显然，$F$ 是连续的。为了证明 $F$ 是一个嵌入，我们只须证明 $F$ 是一个单射（因为 $X$ 是紧致的）。假设 $F(x)=F(y)$。则对于所有的 $i$，有 $\phi_{i}(x)=\phi_{i}(y)$ 和 $h_{i}(x)=h_{i}(y)$。总存在某一个 $i$，使得 $\phi_{i}(x)>0$ [因为 $\sum\phi_{i}(x)=1$]。因此，也必定有 $\phi_{i}(y)>0$，从而 $x,y\in U_{i}$。于是
\begin{align*}
\phi_{i}(x)\cdot g_{i}(x)=h_{i}(x)=h_{i}(y)=\phi_{i}(y)\cdot g_{i}(y).
\end{align*}
又因为 $\phi_{i}(x)=\phi_{i}(y)>0$，所以 $g_{i}(x)=g_{i}(y)$。再根据 $g_{i}\colon U_{i}\to\mathbb{R}^{m}$ 是一个单射，可见 $x=y$。
\end{proof}

在单位分拆的许多应用中，正像上面给出的那样，都只需用到：对于每一个 $x$，$\sum\phi_{i}(x)$ 是正的。然而，在其他一些情形下，就需要一个更强的条件 $\sum\phi_{i}(x)=1$。参见第 50 节。

\subsection*{习题}

\begin{enumerate}
\item 证明每一个流形是正则的，从而是可度量化的。在哪里你要用到 Hausdorff 条件？
\item 设 $X$ 是一个紧致的 Hausdorff 空间。假设对于每一个 $x\in X$，存在 $x$ 的一个邻域 $U$ 和一个正整数 $k$ 使得 $U$ 可以被嵌入到 $\mathbb{R}^{k}$ 中。证明：对于某一正整数 $N$，$X$ 可以嵌入到 $\mathbb{R}^{N}$ 中。
\item 设 $X$ 是一个 Hausdorff 空间，$X$ 的每一点有一个邻域与 $\mathbb{R}^{m}$ 的一个开集同胚。证明：若 $X$ 是紧致的，则 $X$ 是一个 $m$-维流形。
\item $X$ 的子集的一个加标族 $\{A_{\alpha}\}$ 被称为\textbf{点态有限加标族}（point-finite indexed family），如果每一个 $x\in X$，仅仅对于有限多个 $\alpha$ 有 $x\in A_{\alpha}$。

\textbf{引理}（收缩引理）\quad 设 $X$ 是一个正规空间。设 $\{U_{1},U_{2},\cdots\}$ 是 $X$ 的点态有限加标开覆盖。则存在 $X$ 的一个加标开覆盖 $\{V_{1},V_{2},\cdots\}$，使得 $\overline{V}_{n}\subset U_{n}$。
\item 流形定义中的 Hausdorff 条件是关键的，它不能根据定义中的其他条件推出。考虑以下空间：设 $X$ 是 $\mathbb{R}-\{0\}$ 与两点集 $\{p,q\}$ 的并。规定 $X$ 上的拓扑由以下子集组成的基所生成：$\mathbb{R}$ 的所有不包含点 $0$ 的开区间，所有形如 $(-a,0)\cup\{p\}\cup(0,a)$ 的集合以及所有形如 $(-a,0)\cup\{q\}\cup(0,a)$ 的集合，其中 $a>0$。空间 $X$ 称为具有\textbf{双原点的直线}（line with two origins）。
\begin{enumerate}
\item 验证这是某个拓扑的一个基.
\item 证明 $X-\{p\}$ 和 $X-\{q\}$ 分别同胚于 $\mathbb{R}$.
\item 证明 $X$ 满足 $T_{1}$ 公理，但不是一个 Hausdorff 空间.
\item 证明：除了 Hausdorff 条件之外，$X$ 满足 $1$-维流形的所有条件.
\end{enumerate}
\end{enumerate}